% GENERATED FILE. Edit canonical lessons, then run npm run generate:books.
% canonical-source-hash: fnv1a64:75db81572f3ede18
% canonical-lessons: MR-C245-mala2, MR-C245-gotha, MR-C245-masid, MR-C245-mahavidyalay, MR-C245-sangrahalay, MR-R245-first-pass-words-from-chapters-241-242, MR-R245-second-pass-words-from-chapters-243-245

\chapter{Farm plot, Cowshed, Mosque, College, Museum}
\label{ch:mr-farm-plot-cowshed-mosque-college-museum}
\hlchaptermodality{245}

\begin{chapteropening}
\textbf{By the end of this chapter:} \emph{I can say a farm plot, a cowshed, a mosque, a college and a museum.}

Complete the last lesson of chapter 245: I can say a farm plot, a cowshed, a mosque, a college and a museum.
\end{chapteropening}

\section[\texorpdfstring{\mr{मळा} (maḷā)}{मळा (maḷā)}]{\mr{मळा} (maḷā) — a farm plot}

\label{lesson:MR-C245-mala2}



\begin{quote}
Before the new one: say the Marathi for peanuts, then the Marathi for cashews.
\end{quote}

\subsection*{You'll want to know: \mr{मळा}}
\textbf{\mr{मळा}} — \emph{maḷā} — \textquotedblleft{}a farm plot\textquotedblright{}.

\begin{grammarlens}[title={What you've built}]
One more word for this chapter.
\end{grammarlens}

\subsection*{Your turn}
\begin{itemize}
\raggedright
  \item \emph{Say it:} \emph{maḷā}
  \item \emph{Say it:} \emph{maḷā}, once more
  \item \emph{Say it:} say \emph{maḷā}
  \item \emph{Recall:} say the Marathi for peanuts, then the Marathi for cashews, then say \emph{maḷā} again
\end{itemize}

\begin{tcolorbox}[breakable,colback=teal!4,colframe=teal!35!black,title={Before you move on}]
What does \mr{मळा} mean? (\textquotedblleft{}A farm plot\textquotedblright{}.) Say it once more.
\end{tcolorbox}
\section[\texorpdfstring{\mr{गोठा} (goṭhā)}{गोठा (goṭhā)}]{\mr{गोठा} (goṭhā) — a cowshed}

\label{lesson:MR-C245-gotha}



\begin{quote}
Before the new one: say the Marathi for cashews, then the Marathi for a farm plot.
\end{quote}

\subsection*{You'll want to know: \mr{गोठा}}
\textbf{\mr{गोठा}} — \emph{goṭhā} — \textquotedblleft{}a cowshed\textquotedblright{}.

\begin{grammarlens}[title={What you've built}]
One more word for this chapter.
\end{grammarlens}

\subsection*{Your turn}
\begin{itemize}
\raggedright
  \item \emph{Say it:} \emph{goṭhā}
  \item \emph{Say it:} \emph{goṭhā}, once more
  \item \emph{Say it:} say \emph{goṭhā}
  \item \emph{Recall:} say the Marathi for cashews, then the Marathi for a farm plot, then say \emph{goṭhā} again
\end{itemize}

\begin{tcolorbox}[breakable,colback=teal!4,colframe=teal!35!black,title={Before you move on}]
What does \mr{गोठा} mean? (\textquotedblleft{}A cowshed\textquotedblright{}.) Say it once more.
\end{tcolorbox}
\section[\texorpdfstring{\mr{मशीद} (maśīd)}{मशीद (maśīd)}]{\mr{मशीद} (maśīd) — a mosque}

\label{lesson:MR-C245-masid}



\begin{quote}
Before the new one: say the Marathi for a farm plot, then the Marathi for a cowshed.
\end{quote}

\subsection*{You'll want to know: \mr{मशीद}}
\textbf{\mr{मशीद}} — \emph{maśīd} — \textquotedblleft{}a mosque\textquotedblright{}.

\begin{grammarlens}[title={What you've built}]
One more word for this chapter.
\end{grammarlens}

\subsection*{Your turn}
\begin{itemize}
\raggedright
  \item \emph{Say it:} \emph{maśīd}
  \item \emph{Say it:} \emph{maśīd}, once more
  \item \emph{Say it:} say \emph{maśīd}
  \item \emph{Recall:} say the Marathi for a farm plot, then the Marathi for a cowshed, then say \emph{maśīd} again
\end{itemize}

\begin{tcolorbox}[breakable,colback=teal!4,colframe=teal!35!black,title={Before you move on}]
What does \mr{मशीद} mean? (\textquotedblleft{}A mosque\textquotedblright{}.) Say it once more.
\end{tcolorbox}
\section[\texorpdfstring{\mr{महाविद्यालय} (mahāvidyālay)}{महाविद्यालय (mahāvidyālay)}]{\mr{महाविद्यालय} (mahāvidyālay) — a college}

\label{lesson:MR-C245-mahavidyalay}



\begin{quote}
Before the new one: say the Marathi for a cowshed, then the Marathi for a mosque.
\end{quote}

\subsection*{You'll want to know: \mr{महाविद्यालय}}
\textbf{\mr{महाविद्यालय}} — \emph{mahāvidyālay} — \textquotedblleft{}a college\textquotedblright{}.

\begin{grammarlens}[title={What you've built}]
One more word for this chapter.
\end{grammarlens}

\subsection*{Your turn}
\begin{itemize}
\raggedright
  \item \emph{Say it:} \emph{mahāvidyālay}
  \item \emph{Say it:} \emph{mahāvidyālay}, once more
  \item \emph{Say it:} say \emph{mahāvidyālay}
  \item \emph{Recall:} say the Marathi for a cowshed, then the Marathi for a mosque, then say \emph{mahāvidyālay} again
\end{itemize}

\begin{tcolorbox}[breakable,colback=teal!4,colframe=teal!35!black,title={Before you move on}]
What does \mr{महाविद्यालय} mean? (\textquotedblleft{}A college\textquotedblright{}.) Say it once more.
\end{tcolorbox}
\section[\texorpdfstring{\mr{संग्रहालय} (saṅgrahālay)}{संग्रहालय (saṅgrahālay)}]{\mr{संग्रहालय} (saṅgrahālay) — a museum}

\label{lesson:MR-C245-sangrahalay}



\begin{quote}
Before the new one: say the Marathi for a mosque, then the Marathi for a college.
\end{quote}

\subsection*{You'll want to know: \mr{संग्रहालय}}
\textbf{\mr{संग्रहालय}} — \emph{saṅgrahālay} — \textquotedblleft{}a museum\textquotedblright{}.

\begin{grammarlens}[title={What you've built}]
One more word for this chapter.
\end{grammarlens}

\subsection*{Your turn}
\begin{itemize}
\raggedright
  \item \emph{Say it:} \emph{saṅgrahālay}
  \item \emph{Say it:} \emph{saṅgrahālay}, once more
  \item \emph{Say it:} say \emph{saṅgrahālay}
  \item \emph{Recall:} say the Marathi for a mosque, then the Marathi for a college, then say \emph{saṅgrahālay} again
\end{itemize}

\begin{tcolorbox}[breakable,colback=teal!4,colframe=teal!35!black,title={Before you move on}]
What does \mr{संग्रहालय} mean? (\textquotedblleft{}A museum\textquotedblright{}.) Say it once more.
\end{tcolorbox}
\section[(dialogue)]{Words from chapters 241-242 — a first pass}

\label{lesson:MR-R245-first-pass-words-from-chapters-241-242}



\begin{quote}
Say the words for a college and a museum.
\end{quote}

\subsection*{Your turn}
Say these aloud:

\begin{itemize}
\raggedright
  \item a flowerpot, a seedling, a trunk, a lawn, a footpath
  \item leather, metal, a brick, a fine thread, a ladder
\end{itemize}

\begin{tcolorbox}[breakable,colback=teal!4,colframe=teal!35!black,title={Before you move on}]
A flowerpot? (\textbf{\mr{कुंडी}}, \emph{kuṇḍī}.) A seedling? (\textbf{\mr{रोप}}, \emph{rop}.) A trunk? (\textbf{\mr{खोड}}, \emph{khoḍ}.) A lawn? (\textbf{\mr{हिरवळ}}, \emph{hirvaḷ}.) A footpath? (\textbf{\mr{पायवाट}}, \emph{pāyvāṭ}.) Leather? (\textbf{\mr{चामडे}}, \emph{cāmḍe}.) Metal? (\textbf{\mr{धातू}}, \emph{dhātū}.) A brick? (\textbf{\mr{वीट}}, \emph{vīṭ}.) A fine thread? (\textbf{\mr{धागा}}, \emph{dhāgā}.) A ladder? (\textbf{\mr{शिडी}}, \emph{śiḍī}.)
\end{tcolorbox}
\section[(dialogue)]{Words from chapters 243-245 — a second pass}

\label{lesson:MR-R245-second-pass-words-from-chapters-243-245}



\begin{quote}
Say the words for a place and a view.
\end{quote}

\subsection*{Your turn}
Say these aloud:

\begin{itemize}
\raggedright
  \item a place, a view, an outing, a griddle, a wok
  \item mustard seed, coriander leaves, curry leaves, peanuts, cashews
  \item a farm plot, a cowshed, a mosque, a college, a museum
\end{itemize}

\begin{tcolorbox}[breakable,colback=teal!4,colframe=teal!35!black,title={Before you move on}]
A place? (\textbf{\mr{ठिकाण}}, \emph{ṭhikāṇ}.) A view? (\textbf{\mr{देखावा}}, \emph{dekhāvā}.) An outing? (\textbf{\mr{सहल}}, \emph{sahal}.) A griddle? (\textbf{\mr{तवा}}, \emph{tavā}.) A wok? (\textbf{\mr{कढई}}, \emph{kaḍhaī}.) Mustard seed? (\textbf{\mr{मोहरी}}, \emph{moharī}.) Coriander leaves? (\textbf{\mr{कोथिंबीर}}, \emph{kothimbīr}.) Curry leaves? (\textbf{\mr{कढीपत्ता}}, \emph{kaḍhīpattā}.) Peanuts? (\textbf{\mr{शेंगदाणे}}, \emph{śeṅgdāṇe}.) Cashews? (\textbf{\mr{काजू}}, \emph{kājū}.) A farm plot? (\textbf{\mr{मळा}}, \emph{maḷā}.) A cowshed? (\textbf{\mr{गोठा}}, \emph{goṭhā}.) A mosque? (\textbf{\mr{मशीद}}, \emph{maśīd}.) A college? (\textbf{\mr{महाविद्यालय}}, \emph{mahāvidyālay}.) A museum? (\textbf{\mr{संग्रहालय}}, \emph{saṅgrahālay}.)
\end{tcolorbox}
